\clearpage
\hypertarget{native-470-report}{}
\begin{center}
\rule{0.55\linewidth}{0.4pt}\\[2pt]
Astronomische Nachrichten. Band XXI. Nr. 496. Seite 247. 17. Februar 1844.\\[2pt]
\rule{0.55\linewidth}{0.4pt}
\end{center}

Herr Hofrat \textsc{Gauss} hat mir in einem Brief vom 16. die folgenden Beobachtungen des Kometen geschickt:
\begin{center}
\begin{tabular}{l|l|l|l}
1844 Januar 15 & $8^{\rm u}18^{\rm m}17^{\rm s}8$ m. Z. & AR.: $77^\circ20'59''6$ & $\delta$: ---\quad ---\quad --- \\
 & $8\;10\;44.5$\quad --- & AR.: ---\quad ---\quad --- & $\delta$: $+3^\circ49'47''3$ \\
\end{tabular}
\end{center}

Dr. \textsc{Goldschmidt} hatte die Position eines Sterns für diesen Tag zu $78^\circ17'2''3$ $+3^\circ51'2''9$ berechnet. Der Komet ging diesem Stern um $8^{\rm u}18^{\rm m}17^{\rm s}8$ um $3'34''18$ voraus und stand um $8^{\rm u}10^{\rm m}44^{\rm s}5$ mittlerer Zeit $1'15''6$ weiter südlich. Herr Hofrat \textsc{Gauss} hält die Beobachtung für so gut, wie er sie bei der großen Lichtschwäche des Kometen mit dem Ringmikrometer machen konnte, vorausgesetzt selbstverständlich, dass die Position des Sterns gut bestimmt ist.

In einem späteren Brief vom 21. hat Herr Hofrat \textsc{Gauss} mir freundlicherweise den Vergleich mit den zweiten Elementen von \textsc{Goldschmidt} mitgeteilt.
\begin{flushright}\textsc{Schumacher.}\end{flushright}
